%!TEX root = ../main.tex

\chapter{Final Evaluation}\label{ch:final}

The comparisons in Chapters~\ref{ch:rag} and~\ref{ch:kgret} were made on the dev-500
subset, and several design choices were taken by looking at it: the prompt, the
reranker, the settings of entity linking and the choice of graph retriever. Results on
the same data that guided these choices tend to be optimistic, and small differences on
500 questions are hard to separate from noise. The work was therefore completed in two
further steps. First, the systems whose choice was still open were run on the full
validation split (4{,}427 questions), where the remaining decisions were taken and the
final systems were frozen. Second, the frozen systems were run once on the test split of
\gls{mdd} (4{,}335 questions, 241 of them unanswerable), which had not been used before
and was not used for any decision. The generator, the agent-style prompt, the query
rewriter and the evaluation are the same as in the previous chapters.


% ===========================================================================
\section{Decisions on the Validation Split}\label{sec:final:val}

Table~\ref{tab:final:val} compares the candidate systems with the best passage pipeline
(hybrid retrieval, query rewriting and reranking, agent-style prompt) on the full
validation split. The \gls{llm}-only baseline reaches 15.8\% \fone{} and the oracle with
the gold sections 35.5\%.

\begin{table}[ht]
    \centering
    \caption{Full validation split (\%, $n = 4{,}427$). $\Delta$\fone{}: paired bootstrap
    difference to the best passage pipeline with its 95\% interval, on all questions and
    on the answerable ones. Abstain: abstention recall on unanswerable questions / false
    abstentions on answerable ones.}
    \label{tab:final:val}
    \small
    \begin{tabular}{lrrllr}
        \toprule
        System & R@5 & \fone{} & $\Delta$\fone{}, all & $\Delta$\fone{}, answerable & Abstain \\
        \midrule
        Passage pipeline         & 63.7 & 26.9 & ---                       & ---                       & 19.5 / 4.5 \\
        without reranker         & 59.8 & 26.5 & $-0.3$ ($-0.8$, $0.1$)    & $-0.3$ ($-0.7$, $0.0$)    & 19.5 / 4.5 \\
        with the first prompt    & 63.7 & 27.3 & $+0.4$ ($-0.0$, $0.8$)    & $-0.4$ ($-0.7$, $-0.1$)   & 36.0 / 6.0 \\
        A: graph to sections     & 34.0 & 23.0 & $-3.9$ ($-4.5$, $-3.4$)   &                           & 32.3 / 10.4 \\
        B: graph facts           & 36.6 & 24.1 & $-2.7$ ($-3.3$, $-2.1$)   & $-3.4$ ($-3.9$, $-3.0$)   & 30.1 / 7.6 \\
        E: passages + facts      & 63.7 & 27.4 & $+0.5$ ($0.1$, $0.9$)     & $+0.5$ ($0.2$, $0.8$)     & 19.0 / 3.4 \\
        E-random                 & 63.7 & 27.4 & $+0.6$ ($0.2$, $1.0$)     & $-0.0$ ($-0.3$, $0.2$)    & 31.4 / 5.6 \\
        \bottomrule
    \end{tabular}
\end{table}

The reranker raises recall@5 by 3.9 points, while the \fone{} difference is small and not
significant; it was kept. The first, plainer prompt appears 0.4 points better overall,
but only because it abstains much more often (abstention recall 36\% against 20\%), which
is rewarded on the 226 unanswerable questions; on answerable questions it is 0.4 points
worse and gives more false abstentions, so the agent-style prompt was kept. Both
graph-only retrievers are clearly worse than passage retrieval, with facts (B) ahead of
sections (A). Adding graph facts to the passages (E) gains 0.5 points, also on
answerable questions, whereas the random-facts control gains only through extra
abstentions and nothing on answerable questions; E beats it there by 0.6 points (95\%
interval $[0.2, 0.9]$). The systems frozen for the test split are the \gls{llm}-only
baseline, the oracle, the passage pipeline, B, E and E-random.


% ===========================================================================
\section{Results on the Test Split}\label{sec:final:test}

\begin{table}[ht]
    \centering
    \caption{Test split (\%, $n = 4{,}335$). \fone{} with its 95\% bootstrap interval;
    \fone{} on hits and misses for answerable questions with and without a gold span in
    the top-5.}
    \label{tab:final:test}
    \small
    \begin{tabular}{lrlrrr}
        \toprule
        System & R@5 & \fone{} & Hit / miss & Abstain rec. & False abst. \\
        \midrule
        \gls{llm} only          & ---  & 15.1 ($14.7$, $15.5$) & ---           & 14.5 & 11.8 \\
        Oracle                  & 100  & 34.2 ($33.6$, $34.9$) & 35.1 / ---    & 18.7 &  2.4 \\
        Passage pipeline        & 66.6 & 26.9 ($26.3$, $27.5$) & 31.9 / 17.4   & 22.4 &  5.2 \\
        B: graph facts          & 38.7 & 23.8 ($23.2$, $24.4$) & 28.6 / 19.9   & 32.0 & 10.0 \\
        E: passages + facts     & 66.6 & 27.5 ($27.0$, $28.2$) & 32.4 / 19.0   & 19.1 &  3.3 \\
        E-random                & 66.6 & 27.3 ($26.7$, $27.9$) & 31.8 / 17.6   & 29.9 &  6.5 \\
        \bottomrule
    \end{tabular}
\end{table}

\begin{table}[ht]
    \centering
    \caption{Test split: $\Delta$\fone{} against the passage pipeline (points, paired
    bootstrap, 95\% interval).}
    \label{tab:final:delta}
    \small
    \begin{tabular}{llll}
        \toprule
        & All & Answerable & Passages miss gold \\
        \midrule
        B: graph facts      & $-3.1$ ($-3.7$, $-2.5$) & $-3.8$ ($-4.3$, $-3.3$) & $+1.0$ ($0.4$, $1.8$) \\
        E: passages + facts & $+0.6$ ($0.2$, $1.1$)   & $+0.8$ ($0.5$, $1.2$)   & $+1.6$ ($1.1$, $2.1$) \\
        E-random            & $+0.4$ ($-0.0$, $0.8$)  & $-0.0$ ($-0.3$, $0.3$)  & $+0.2$ ($-0.1$, $0.6$) \\
        E vs E-random       & $+0.2$ ($-0.2$, $0.7$)  & $+0.9$ ($0.6$, $1.2$)   & $+1.4$ ($0.9$, $1.9$) \\
        \bottomrule
    \end{tabular}
\end{table}

The test results (Tables~\ref{tab:final:test} and~\ref{tab:final:delta}) reproduce the
validation results almost exactly. The ordering of the systems is the same: the
\gls{llm} alone reaches 15.1\% \fone{}, retrieval over the graph alone 23.8\%, passage
retrieval 26.9\%, passages with graph facts 27.5\% and the oracle 34.2\%. Retrieval over
the graph alone is 3.1 points worse than passage retrieval, and 3.8 points on answerable
questions. Adding graph facts to the retrieved passages gains 0.6 points overall and 0.8
on answerable questions, both significant, and 1.6 points on the questions where passage
retrieval misses the gold section. On the 2{,}206 answerable questions where at least one
graph fact comes from the gold section the gain is 1.4 points (interval $[1.0, 1.8]$),
while on the others it is not significant. Random facts gain nothing on answerable
questions and are beaten by the graph facts by 0.9 points; their small overall gain comes
again from more abstentions (abstention recall 30\% against 22\%), which is rewarded on
unanswerable questions but also produces more false abstentions.

The effect of the graph facts, which was only borderline on dev-500, is therefore
confirmed on two large sets that played no role in its design, and the random-facts
control shows that it comes from the content of the facts rather than from the longer
context. It is small: about half a point of \fone{}, against a gap of 7.3 points between
the passage pipeline and the oracle. On the test split, a gold section is among the top
five retrieved sections for 67\% of the answerable questions, with an \fone{} of 31.9\%
on these against 17.4\% on the others, so retrieval remains the main limitation.


% ===========================================================================
\section{Notes on the Runs}\label{sec:final:notes}

The provider's content filter blocked 8 of about 66{,}000 generated answers (7 for the
first prompt on the validation split and 1 for the \gls{llm}-only baseline on the test
split); these questions are left out of the scores of the two systems. During the test
run a consistency check, which verifies that the passages in variant E are exactly those
of the passage pipeline, revealed a race condition: when two questions led to the same
rewriting prompt, both were sent to the model at the same time and could receive
different rewrites, of which only one was cached. Seven test questions were affected in
the retrieval data of the graph pipelines. Identical prompts are now sent only once, the
affected data were rebuilt and the graph pipelines rerun; the validation split and
dev-500 passed the same check.
